\subsection{Complete decision traces and finite-pool risk}
\label{sec:7-budget-trace}

Replaying the recorded calibration choices confirms the finite-pool
calculation: all 2,304 sampled decisions and twelve full-pool decisions
match enumeration. Their recorded discrepancy curves are monotone, and
summing per-query counts reproduces all 18,432 tested grid counts
(\cref{eq:6-finite-pool}). This validates the observed traces; it does
not establish nesting for every distance-tied subset or query ordering
(\cref{sec:6-3}).

At $S=32$, maximum individual loss reaches 0.0294, versus at most
0.0006 for full-pool calibration.

For arXiv at 128 probes, the predicted and observed probabilities of
exceeding the loss target are 0.965 versus 0.984 with $S=32$,
0.875 versus 0.875 with $S=64$, and 0.620 versus 0.625 with $S=128$
(\cref{fig:7-finite-pool}). The largest of 36 gaps occurs on BGE at
512 probes: 0.519 predicted versus 0.375 observed, a difference of
0.144. These predictions concern samples from the fixed calibration
pool, using each selected budget's measured held-out quality; calibration
itself uses no labels.

\begin{figure}[!htbp]
\centering
\includegraphics[width=\linewidth]{supplement/finite-pool.pdf}
\caption[Finite-pool risk prediction and observation]{Finite-pool prediction and observed arXiv loss crossings. Bars
show one estimated standard error over 64 independently drawn subsets
conditional on the fixed pool, not uncertainty over future workloads.}
\label{fig:7-finite-pool}
\end{figure}

Ten alternating timing repeats show a low-budget plateau, rather than
a general rule that each smaller budget is faster.

\subsection{Direct calibration-plus-reuse measurement}
\label{sec:7-stream-cost}

Four settings compare fixed budgets 64 and 100,
128-query calibration and full-pool calibration across five rotated
rounds of 500 reuse batches. Every batch repeats the same held-out query
pool. This measures amortisation on a stable workload, not performance
under new-query arrivals or drift. Timing includes sample gathering,
calibration, budget installation and initial graph setup, but excludes
construction, warm-up and scan-scheduler calibration.

\par\medskip\noindent\begin{minipage}{\linewidth}
\centering\captionsetup{hypcap=false}
\captionof{table}[Direct costs for 500 reuse batches]{Direct total milliseconds for 500 reuse batches, averaged over
five rounds, including budget calibration for sampled and FULL policies.
Each batch has 500 queries. Quality is discussed with the timing results.}
\label{tab:7-stream-cost}
\footnotesize\setlength{\tabcolsep}{4pt}
\begin{tabular}{@{}lrrrr@{}}
\toprule
Dataset/probes & Fixed 100 & Fixed 64 & Rule $S=128$ & FULL \\
\midrule
arXiv/128 & 1224 & 1102 & 1291 & 1680 \\
OpenAI-1536/128 & 1491 & 1264 & 852 & 860 \\
OpenAI-3072/128 & 2757 & 2160 & 1769 & 1783 \\
BGE/512 & 5185 & 5018 & 4921 & 5062 \\
\bottomrule
\end{tabular}

\end{minipage}\par\medskip

Both calibrated OpenAI-3072 policies choose budget 4 with recall 0.9252,
matching fixed 100. OpenAI-1536 full-pool calibration and fixed 100 give
0.9280, while sampling gives 0.9276. BGE at 512 probes gives 0.9658
for sampling, 0.9670 for fixed 100 and 0.9664 for full-pool calibration.
ArXiv's full-pool policy selects the slower budget 200 to meet quality,
while its fixed controls fail the target.

For 128-query calibration, mean cumulative cost first falls below
fixed 100 after four batches on each OpenAI dataset. Against fixed 64,
it crosses after six batches on OpenAI-1536 and nine on OpenAI-3072.
Full-pool calibration requires nine and fourteen batches on OpenAI-1536,
and eleven and twenty-seven on OpenAI-3072, respectively
(\cref{fig:7-stream-cost}). These measured crossings complement the
steady-rate model; policies with different recall cannot be ranked by
time alone.

\par\medskip\noindent\begin{minipage}{\linewidth}
\captionsetup{hypcap=false}
\centering
\includegraphics[width=\linewidth]{supplement/stream-cost.pdf}
\captionof{figure}[Direct cumulative calibrated-reuse time]{Direct mean cumulative time, including budget calibration and
first-batch setup. Each curve replays one fixed held-out query pool.
Policies can have different recall; timings alone are not a quality-matched
ranking. Dataset/probe settings match \cref{tab:7-stream-cost}.}
\label{fig:7-stream-cost}
\end{minipage}\par\medskip
